\documentclass[11pt]{article}
\usepackage[margin=1in]{geometry}
\usepackage[T1]{fontenc}
\usepackage[utf8]{inputenc}
\usepackage[expansion=false]{microtype}
\usepackage{amsmath,amssymb,booktabs,longtable,array,siunitx,enumitem,hyperref,xurl}
\hypersetup{colorlinks=true,linkcolor=blue,urlcolor=blue}
\setlist{nosep}
\setlength{\parindent}{0pt}
\setlength{\parskip}{0.4em}
\setlength{\LTleft}{0pt}
\setlength{\LTright}{0pt}
\setlength{\tabcolsep}{4pt}
\setlength{\emergencystretch}{2em}
\allowdisplaybreaks
\newcolumntype{L}[1]{>{\raggedright\arraybackslash}p{#1}}
\newcommand{\srcref}[1]{\nolinkurl{#1}}
\newcommand{\filepath}[1]{\nolinkurl{#1}}
\newcommand{\code}[1]{\nolinkurl{#1}}

\title{Physics and Mathematics Reference\\Recreated Heat Transfer Study Model}
\date{Generated 2026-04-03}

\begin{document}
\maketitle

\section{Purpose and scope}
This document describes the physics and mathematics implemented in the recreated study script
\srcref{heat_transfer_study_recreated.py}, including:
\begin{itemize}[leftmargin=*]
\item legacy winter heating from the exported heater/aeration sweep;
\item simplified winter wall-blanket heating;
\item legacy summer spot-cooler plus assist-blower cooling;
\item simplified summer single vertical frozen-bottle column cooling;
\item shared pressure-drop, moisture-loss, two-node bed, radial plume, and year-round layers.
\end{itemize}

Primary implementation anchors:
\begin{itemize}[leftmargin=*]
\item runtime dispatch and mode selection: \srcref{heat_transfer_study_recreated.py:L9427-L9557};
\item heating mode config: \srcref{heat_transfer_study_recreated.py:L6500-L6520};
\item cooling mode config: \srcref{heat_transfer_study_recreated.py:L6521-L6633};
\item optimization config and temperature-goal toggles: \srcref{heat_transfer_study_recreated.py:L2506-L2586}.
\end{itemize}

\section{Execution modes and configuration controls}
Two run layers can be enabled independently via \code{run\_modes.heating} and \code{run\_modes.cooling}.
The recreated profiles are:
\begin{itemize}[leftmargin=*]
\item legacy+ profile:
\srcref{study_config_legacy_plus.json:L8-L399};
\item simplified vertical/wall-blanket profile:
\srcref{study_config_simplified_vertical.json:L8-L446}.
\end{itemize}

Mode-selecting fields:
\begin{itemize}[leftmargin=*]
\item \code{heating\_mode.method}:
\code{legacy\_export\_from\_matlab} or \code{wall\_blanket\_conduction}
(\srcref{study_config_simplified_vertical.json:L8-L23});
\item \code{cooling\_mode.method}:
\code{spot\_cooler\_with\_assist\_blower} or
\code{vertical\_frozen\_bottle\_column}
(\srcref{study_config_simplified_vertical.json:L236-L239}).
\end{itemize}

Optimization and seasonal-goal toggles:
\begin{itemize}[leftmargin=*]
\item \code{optimization.method} and \code{cooling\_mode.optimization.method}:
\code{nsga2} or \code{differential\_evolution};
\item \code{temperature\_goal\_mode}:
\code{legacy\_extreme\_node} or \code{seasonal\_threshold};
\item threshold targets:
\code{winter\_minimum\_target\_C}, \code{summer\_maximum\_target\_C}.
\end{itemize}
Config anchors:
\srcref{study_config_legacy_plus.json:L264-L330},
\srcref{study_config_simplified_vertical.json:L311-L377}.

\section{Thermophysical and psychrometric closures}
Air properties (ideal-gas and Sutherland-law transport):
\begin{align}
\rho &= \frac{p}{R_a T}, \\
\mu &= \mu_0\left(\frac{T}{T_0}\right)^{3/2}\frac{T_0+S}{T+S}, \\
k &= \frac{\mu c_p}{Pr}.
\end{align}
Implemented in \srcref{heat_transfer_study_recreated.py:L4972-L4980}.

Humidity-ratio relations:
\begin{align}
w &= 0.62198\,\frac{p_v}{p-p_v}, \\
p_v &= \phi\,p_{\mathrm{sat}}(T).
\end{align}
Implemented in
\srcref{heat_transfer_study_recreated.py:L5016-L5029},
with dew-point inversion in
\srcref{heat_transfer_study_recreated.py:L5047-L5060}.
Water saturation pressure and latent heat use the imported IAPWS backend
(\srcref{heat_transfer_study_recreated.py:L4983-L5014}).

\section{Convection and friction correlations}
Wire-side external forced convection uses Churchill--Bernstein:
\begin{equation}
Nu = 0.3 + \frac{0.62 Re^{1/2} Pr^{1/3}}{\left(1+(0.4/Pr)^{2/3}\right)^{1/4}}
\left(1+\left(\frac{Re}{282000}\right)^{5/8}\right)^{4/5}.
\end{equation}
Implemented in \srcref{heat_transfer_study_recreated.py:L5150-L5157}.

Tube friction factor uses Churchill's smooth-pipe form:
\begin{equation}
f = 8\left[\left(\frac{8}{Re}\right)^{12} + \frac{1}{(A+B)^{3/2}}\right]^{1/12},
\end{equation}
with laminar fallback \(f=64/Re\), implemented in
\srcref{heat_transfer_study_recreated.py:L5159-L5166}.

Internal tube Nusselt closure:
\begin{itemize}[leftmargin=*]
\item laminar developing flow: \(Nu = 3.66 + \frac{0.0668\,Gz}{1+0.04\,Gz^{2/3}}\),
\(Gz=Re\,Pr\,D/L\);
\item turbulent flow: Gnielinski-type relation using \(f\).
\end{itemize}
Implemented in \srcref{heat_transfer_study_recreated.py:L5168-L5177}.

Natural convection for wall-loss closure:
\begin{itemize}[leftmargin=*]
\item vertical surfaces: Churchill--Chu style \(Nu(Ra,Pr)\);
\item horizontal surfaces: stable/unstable orientation piecewise \(Nu(Ra)\).
\end{itemize}
Implemented in \srcref{heat_transfer_study_recreated.py:L5062-L5099}.

\section{Overall U and wall-loss resistances}
The cylindrical overall \(U\) on the outer-area basis combines inner convection,
wall conduction, optional insulation conduction, and outer convection:
\begin{equation}
\frac{1}{U_o} =
\frac{D_o}{D_i h_i}
 + \frac{D_o \ln(D_w/D_i)}{2k_w}
 + \frac{D_o \ln(D_o/D_w)}{2k_{ins}}
 + \frac{1}{h_o}.
\end{equation}
Implemented in \srcref{heat_transfer_study_recreated.py:L5179-L5215}.

\section{Legacy winter heating path (exported heater/aeration design points)}
For legacy heating, each sampled design point carries electrical topology,
heater-tube segment results, aeration-tube release, pressure loss, and bed-node
temperatures via \code{simulate\_heating\_reference\_point}
(\srcref{heat_transfer_study_recreated.py:L5521-L6283}).

Electrical grouping uses strings of series tubes under one supply voltage.
For one string:
\begin{align}
I_{\mathrm{string}} &= \frac{V}{n_s R_{\mathrm{tube}}}, \\
P_{\mathrm{tube}} &= I_{\mathrm{string}}^2 R_{\mathrm{tube}}.
\end{align}
The full point aggregates all strings/tubes.

Heater-tube segment energy update:
\begin{align}
\dot q_{\mathrm{gen},i} &= P_{\mathrm{tube}}\frac{\Delta x}{L}, \\
\dot q_{\mathrm{loss},i} &= U_{\mathrm{loss}} P_{\mathrm{exp}}\Delta x
\max(T_{air,i}-T_\infty,0), \\
T_{air,i+1} &= T_{air,i} +
\frac{\dot q_{\mathrm{gen},i}-\dot q_{\mathrm{loss},i}}
{\dot m c_p}.
\end{align}
Implemented in \srcref{heat_transfer_study_recreated.py:L5248-L5347}.

\section{Simplified winter wall-blanket conduction path}
When \code{heating\_mode.method = wall\_blanket\_conduction}, electrical aeration
heating is replaced by wall blankets under insulation
(\srcref{heat_transfer_study_recreated.py:L5523-L5630}).
No aeration-driven heating or nichrome heater-tube heating is used in this
winter branch.

Configured blanket hardware basis (per blanket):
\begin{itemize}[leftmargin=*]
\item dimensions \(0.250\times1.740\times0.0015\;\si{m}\) (10 in \(\times\) 70 in \(\times\) 1.5 mm);
\item rated electrical power \(1200\;\si{W}\) at \(110\;\si{V}\);
\item modeled max blanket surface temperature \(120\,^\circ\mathrm{C}\).
\end{itemize}

Blanket surface temperature is duty-ramped:
\begin{equation}
T_{blanket} = T_\infty + \delta\,(T_{blanket,\max}-T_\infty), \qquad 0\le\delta\le 1.
\end{equation}

Inward and outward heat fluxes:
\begin{align}
q''_{\mathrm{in}} &= \frac{\max(T_{blanket}-T_{bed},0)}{R_{\mathrm{in}}}, \\
q''_{\mathrm{out}} &= \frac{\max(T_{blanket}-T_\infty,0)}{R_{\mathrm{out}}}.
\end{align}

Total powers:
\begin{align}
Q_{\mathrm{to\,bed}} &= N_b A_b q''_{\mathrm{in}}\,\eta_{coupling}, \\
Q_{\mathrm{to\,amb}} &= N_b A_b q''_{\mathrm{out}}, \\
P_{\mathrm{elec}} &= Q_{\mathrm{to\,bed}} + Q_{\mathrm{to\,amb}}, \\
I_{\mathrm{total}} &= \frac{P_{\mathrm{elec}}}{V_{supply}}.
\end{align}

Thermal response still uses the same two-node bed network and the same wall-loss
model as all other branches.

\section{Summer inlet-conditioning models}
\subsection{Spot-cooler plus assist-blower mode}
For \code{spot\_cooler\_with\_assist\_blower}, requested sensible cooling is
\begin{equation}
Q_{\mathrm{req}} = \dot m c_p (T_{amb}-T_{target})_+.
\end{equation}
Actual load is capped by configured/rated sensible capacity, then
\begin{equation}
T_{supply} = T_{amb} - \frac{Q_{\mathrm{actual}}}{\dot m c_p}.
\end{equation}
Condensate is estimated from psychrometric carrying-capacity difference with an
optional rated dehumidification cap. Implemented in
\srcref{heat_transfer_study_recreated.py:L3512-L3598}.

\subsection{Simplified vertical frozen-bottle column mode}
For \code{vertical\_frozen\_bottle\_column}, a single central vertical tube is
used; fan flow is downward through packed frozen bottles
(\srcref{heat_transfer_study_recreated.py:L3601-L3699},
\srcref{heat_transfer_study_recreated.py:L3703-L4108}).

Column superficial velocity:
\begin{equation}
u_s = \frac{\dot V}{A_{col}}, \qquad
A_{col} = \frac{\pi D_{col}^2}{4}.
\end{equation}

Convective coefficient and effective UA:
\begin{align}
h_{col} &= h_0 + C_v u_s^{0.8}, \\
UA_{col} &= h_{col} A_{exp}
\quad \text{(or fixed override if configured)}.
\end{align}

Unconstrained outlet from an exponential UA model:
\begin{equation}
T_{out,ideal} = T_{sink} + (T_{amb}-T_{sink})
\exp\!\left(-\frac{UA_{col}}{\dot m c_p}\right).
\end{equation}

Frozen-bottle capacity cap:
\begin{align}
E_{ice} &= m_w\!\left[c_{p,ice}(0-T_{ice,0}) + h_{fus}
+ c_{p,liq}(T_{post}-0)\right], \\
\dot Q_{cap} &= \frac{E_{ice}}{\tau_{replace}}, \\
\dot Q_{actual} &= \min(\dot Q_{ideal},\dot Q_{cap}).
\end{align}

The resulting \(T_{supply}\), humidity, and condensate then feed the same
aeration/bed solvers used by the other summer branch.

\section{Aeration pressure model and closed-end perforated-tube release}
Both heating and cooling branches use the same core release and pressure closures.

Per-segment perforation release:
\begin{equation}
\dot m_{\mathrm{rel},i}
= C_d A_{h,i}\sqrt{2\rho_i\,(p_{tube,i}-p_{bed,i})_+}.
\end{equation}
Implemented in
\srcref{heat_transfer_study_recreated.py:L3154-L3214},
\srcref{heat_transfer_study_recreated.py:L3366-L3436}.

Segmentwise bed backpressure uses an Ergun pressure gradient:
\begin{equation}
\frac{d p}{d x}
=
150\frac{(1-\epsilon)^2}{\epsilon^3}
\frac{\mu u_s}{d_p^2}
+
1.75\frac{1-\epsilon}{\epsilon^3}
\frac{\rho u_s^2}{d_p},
\end{equation}
with local path length to open top or vent.
Implemented in
\srcref{heat_transfer_study_recreated.py:L3216-L3311},
\srcref{heat_transfer_study_recreated.py:L3334-L3364}.

Closed-end pressure is solved by scalar shooting:
\begin{align}
\dot m_{\mathrm{rem},i+1} &= \dot m_{\mathrm{rem},i} - \dot m_{\mathrm{rel},i}, \\
p_{i+1} &= \max(p_i-\Delta p_{fric,i},\,p_{bed,i}).
\end{align}
The inlet gauge pressure is bisected until terminal residual flow is below
tolerance. Implemented in
\srcref{heat_transfer_study_recreated.py:L3442-L3510}.

\section{Distribution-network losses}
Legacy manifold-network losses include:
\begin{itemize}[leftmargin=*]
\item header entry and header-to-splitter contraction;
\item splitter body/tee/valve surrogate losses (Hooper 2-K style);
\item branch connector friction;
\item connector-to-branch contraction/expansion losses.
\end{itemize}
Implemented in
\srcref{heat_transfer_study_recreated.py:L5787-L6323}.

In simplified vertical-column summer mode, these can be bypassed via
\code{vertical\_column.bypassDistributionNetwork = true}
(\srcref{heat_transfer_study_recreated.py:L3801-L3817}).

\section{Aeration segment energy, jet transfer, and evaporation}
Each aeration segment contributes:
\begin{align}
Q_{\mathrm{to\,bed},i}
&=
Q_{\mathrm{wall},i}
+
Q_{\mathrm{jet},i}
-
Q_{\mathrm{evap},i}, \\
Q_{\mathrm{jet},i}
&= \dot m_{\mathrm{rel},i} c_p (T_{\mathrm{rel},i}-T_{bed}).
\end{align}

Evaporation rate is triple-limited:
\begin{equation}
\dot m_{evap,i}
=
\min\!\left(
\dot m_{diff,i},
\dot m_{capacity,i},
\dot m_{heat,i}
\right),
\qquad
Q_{\mathrm{evap},i}=\dot m_{evap,i} h_{fg}.
\end{equation}
Implemented in
\srcref{heat_transfer_study_recreated.py:L3160-L3200},
\srcref{heat_transfer_study_recreated.py:L5406-L5486},
\srcref{heat_transfer_study_recreated.py:L3892-L3999}.

\section{Bin wall losses and two-node bed closure}
The bin loss model assembles side, bottom, and top UA based on top state
(open, covered+vent, covered) and bottom boundary mode:
\srcref{heat_transfer_study_recreated.py:L6324-L6464}.

Lumped requirement:
\begin{equation}
Q_{req} = UA_{tot}(T_{set}-T_{amb}).
\end{equation}
Implemented in \srcref{heat_transfer_study_recreated.py:L6466-L6486}.

Two-node steady state:
\begin{equation}
\begin{bmatrix}
UA_b + UA_{12} & -UA_{12} \\
-UA_{12} & UA_t + UA_{12}
\end{bmatrix}
\begin{bmatrix}T_b\\T_t\end{bmatrix}
=
\begin{bmatrix}
Q_b + UA_b T_{amb}\\
Q_t + UA_t T_{amb}
\end{bmatrix},
\end{equation}
with \(Q_b=f_b Q_{bed}\), \(Q_t=(1-f_b)Q_{bed}\).
Implemented in \srcref{heat_transfer_study_recreated.py:L6488-L6498}.

\section{Radial plume and habitat exclusion post-processing}
The post-processing plume uses release-air cooling away from the tube with
configurable linear plume-radius spread:
\begin{equation}
r_p(x)=\min\!\left(r_{max},\,r_0+\alpha x\right).
\end{equation}

Sensible and latent sink terms are applied along radius; the ``safe distance''
is the first radius where \(T(x)\le T_{threshold}\).
Implemented in
\srcref{heat_transfer_study_recreated.py:L1365-L1708}.

Tube-layout and exclusion-area mapping is then computed on a 2-D cross-section
grid using either explicit centers, custom single row, or rule-based patterns:
\srcref{heat_transfer_study_recreated.py:L1709-L2050}.

\section{Constraint residuals and objective metrics}
\subsection{Heating metrics}
Heating residuals include current, wire temperature, outlet-air temperature,
pressure drop, moisture loss, hole velocity, bottom/top min and max limits,
spread limit, and heat shortfall:
\begin{equation}
r_k(x)=\max(0,g_k(x)).
\end{equation}
Implemented in \srcref{heat_transfer_study_recreated.py:L2610-L2724}.

Legacy heating objective family:
\begin{itemize}[leftmargin=*]
\item \code{power\_W};
\item \code{minus\_min\_bed\_temp\_C};
\item \code{temp\_spread\_C}.
\end{itemize}

Threshold heating objective family:
\begin{itemize}[leftmargin=*]
\item \code{winter\_threshold\_shortfall\_C};
\item \code{power\_W};
\item \code{temp\_spread\_C}.
\end{itemize}

\subsection{Cooling metrics}
Cooling residuals include assist-blower margin, pressure drop, moisture loss,
hole velocity, temperature bounds/spread, and nonpositive cooling delivery.
Implemented in \srcref{heat_transfer_study_recreated.py:L4112-L4225}.

Legacy cooling objective family:
\begin{itemize}[leftmargin=*]
\item \code{max\_bed\_temp\_C};
\item \code{mean\_bed\_temp\_C};
\item \code{combined\_cooling\_power\_W}.
\end{itemize}

Threshold cooling objective family:
\begin{itemize}[leftmargin=*]
\item \code{summer\_threshold\_excess\_C};
\item \code{combined\_cooling\_power\_W};
\item \code{totalFlow\_Lpm}.
\end{itemize}

\section{Optimization methods}
\subsection{NSGA-II mode}
When \code{method = nsga2}, constrained NSGA-II is applied directly on the
candidate design-point set:
\srcref{heat_transfer_study_recreated.py:L2889-L2919},
\srcref{heat_transfer_study_recreated.py:L4317-L4347}.

\subsection{Differential-evolution mode}
When \code{method = differential\_evolution}, the solver runs multiple weighted
scalarizations and then keeps the feasible non-dominated set.

For one weight vector \(w_j\):
\begin{equation}
J_j(x)
=
\sum_i w_{j,i}\,
\frac{f_i(x)-f_{i,\min}}{\Delta f_i}
\;+\;
\lambda \sum_k r_k(x),
\end{equation}
where \(r_k(x)\) are nonnegative constraint residuals.

Heating DE implementation:
\srcref{heat_transfer_study_recreated.py:L2921-L3001}.
Cooling DE implementation:
\srcref{heat_transfer_study_recreated.py:L4349-L4424}.

\section{Seasonal threshold objective mode}
With \code{temperature\_goal\_mode = seasonal\_threshold}:
\begin{itemize}[leftmargin=*]
\item winter heating objective/priority includes
\code{winter\_threshold\_shortfall\_C} relative to
\code{winter\_minimum\_target\_C};
\item summer cooling objective/priority includes
\code{summer\_threshold\_excess\_C} relative to
\code{summer\_maximum\_target\_C}.
\end{itemize}
Implementation anchors:
\srcref{heat_transfer_study_recreated.py:L2556-L2586},
\srcref{heat_transfer_study_recreated.py:L2633-L2645},
\srcref{heat_transfer_study_recreated.py:L4138-L4151}.

\section{Year-round layer}
The annual model builds a representative climate year via periodic Gaussian-kernel
regression from monthly means/standard deviations, then applies freeze-day and
heat-wave overrides. Heat-wave spell starts are selected by Gaussian-centered
repulsion scoring, and spell anomaly magnitudes are sampled from a seeded
truncated normal bounded by NOAA-derived monthly spell min/max anomaly arrays
(derived-file source:
\code{Heat Transfer Study/data/climate/noaa\_heatwave\_spell\_anomaly\_bounds\_USW00014740\_2016\_2025.json}):
\srcref{heat_transfer_study_recreated.py:L6910-L7023}.

Daily simulation:
\begin{enumerate}[leftmargin=*]
\item compute daily required heating/cooling from the same lumped UA model;
\item re-simulate selected heating or cooling reference point at that day's ambient;
\item apply duty, unmet-load accounting, electricity use, and water use.
\end{enumerate}
Implemented in
\srcref{heat_transfer_study_recreated.py:L7025-L7133} and
\srcref{heat_transfer_study_recreated.py:L7134-L7265}.

\section{Symbol glossary}
\begin{longtable}{L{0.22\textwidth}L{0.57\textwidth}L{0.17\textwidth}}
\toprule
Symbol & Meaning & Units \\
\midrule
\endfirsthead
\toprule
Symbol & Meaning & Units \\
\midrule
\endhead
\(Q_{bed}\) & Net heat delivered to bed (positive heating, negative cooling) & W \\
\(Q_{req}\) & Lumped required load at setpoint and ambient & W \\
\(UA\) & Overall heat-loss conductance & W/K \\
\(T_b, T_t\) & Bottom and top bed-node temperatures & \(^{\circ}\)C \\
\(\dot m_{\mathrm{rel},i}\) & Segment release mass flow from perforations & kg/s \\
\(\Delta p\) & Pressure drop / required static pressure & Pa \\
\(\epsilon\) & Bed porosity in Ergun closure & -- \\
\(d_p\) & Representative bed particle diameter & m \\
\(h_{fg}\) & Latent heat of vaporization & J/kg \\
\(\dot m_{evap}\) & Evaporation mass flow & kg/s \\
\(u_s\) & Superficial velocity (context-dependent) & m/s \\
\(\delta\) & Wall-blanket duty fraction & -- \\
\bottomrule
\end{longtable}

\section{Implementation files}
\begin{itemize}[leftmargin=*]
\item model script:
\filepath{S:\textbackslash AcademicResources\textbackslash UCONN\textbackslash 10. Spring 2026\textbackslash CHEG 4143W\textbackslash Project\textbackslash VermiHeating\textbackslash Heat Transfer Study\textbackslash Reduced\textbackslash heat\_transfer\_study\_recreated.py}
\item legacy-plus config:
\filepath{S:\textbackslash AcademicResources\textbackslash UCONN\textbackslash 10. Spring 2026\textbackslash CHEG 4143W\textbackslash Project\textbackslash VermiHeating\textbackslash Heat Transfer Study\textbackslash Reduced\textbackslash study\_config\_legacy\_plus.json}
\item simplified config:
\filepath{S:\textbackslash AcademicResources\textbackslash UCONN\textbackslash 10. Spring 2026\textbackslash CHEG 4143W\textbackslash Project\textbackslash VermiHeating\textbackslash Heat Transfer Study\textbackslash Reduced\textbackslash study\_config\_simplified\_vertical.json}
\end{itemize}

\end{document}
